% Problem record op_7e7e4a25fef5c994. This TeX file is derived from
% database/problems_json/op_7e7e4a25fef5c994.json by scripts/sync-tex.mjs; edit the JSON record
% and rerun the script rather than editing this file.
\section{Transpose degradability beyond degradability}
\label{sec:op_7e7e4a25fef5c994}

\paragraph{Problem.}
Does there exist a finite-dimensional transpose-degradable quantum channel
that is not degradable?  Let $V:A\to B\otimes E$ be an isometry defining a
channel and a complementary channel by
\begin{equation}
  \Phi(X):=\operatorname{Tr}_E(VXV^\dagger),
  \qquad
  \Phi^c(X):=\operatorname{Tr}_B(VXV^\dagger).
  \label{eq:p53-complementary-pair}
\end{equation}
Equation~\eqref{eq:p53-complementary-pair} fixes the output space $B$ and
environment space $E$.  For the transpose $\mathsf T_E$ in a fixed basis of
$E$, transpose degradability means that a completely positive trace-preserving
map $\mathcal D:\mathcal L(B)\to\mathcal L(E)$ satisfies
\begin{equation}
  \mathsf T_E\circ\Phi^c=\mathcal D\circ\Phi.
  \label{eq:p53-transpose-degradable}
\end{equation}
Ordinary degradability instead requires a completely positive
trace-preserving map
$\widetilde{\mathcal D}:\mathcal L(B)\to\mathcal L(E)$ satisfying
\begin{equation}
  \Phi^c=\widetilde{\mathcal D}\circ\Phi.
  \label{eq:p53-degradable}
\end{equation}
The question is whether Eq.~\eqref{eq:p53-transpose-degradable} can hold while
no map satisfying Eq.~\eqref{eq:p53-degradable} exists.

\subsection*{Status}

Solved

\subsection*{Source}

Singh and Datta explicitly ask whether transpose-degradable channels differ
from ordinary degradable channels
\sourcecite{ref:p53-singh-datta}{SD22}.  Br{\'a}dler had posed the same
strict-separation question using the earlier \lq\lq conjugate degradable\rq\rq{}
terminology \sourcecite{ref:p53-bradler}{Bra15}.

\subsection*{Progress}

\begin{itemize}
  \item Writing $J(\mathcal N)$ for the Choi operator of a channel
  $\mathcal N$, Eq.~\eqref{eq:p53-transpose-degradable} implies
  \begin{equation}
    J(\Phi^c)^{T_E}\geq0,
    \qquad
    Q(\Phi)=Q^{(1)}(\Phi)
    :=\max_{\rho_A}
       \left[S(\Phi(\rho_A))-S(\Phi^c(\rho_A))\right].
    \label{eq:p53-ppt-and-capacity}
  \end{equation}
  Thus Eq.~\eqref{eq:p53-ppt-and-capacity} gives a PPT complementary Choi
  operator and additive coherent information, but neither property supplies an
  ordinary degrading map \sourcecite{ref:p53-singh-datta}{SD22}.

  \item If $J(\Phi^c)$ is separable, then $\Phi^c$ is entanglement breaking and
  $\Phi$ is a Hadamard channel, hence degradable.  Any strict example must
  therefore have a PPT-entangled complementary Choi operator.  Moreover,
  transpose-degradable pcubed channels are always ordinarily degradable, so
  that structured family contains no strict example
  \sourcecite{ref:p53-bradler}{Bra15},
  \sourcecite{ref:p53-siddhu-griffiths}{SG16}.

  \item The universal-cloning family also provides no strict example.  For all
  $d\geq2$ and $N,K\geq1$, the optimal symmetric cloner
  $\mathcal C_{N\to N+K}^{(d)}$ and optimal pure-state transposition channel
  $\mathcal T_{N\to K}^{(d)}$ obey
  \begin{equation}
    \left(\mathcal C_{N\to N+K}^{(d)}\right)^c
      =\mathcal T_{N\to K}^{(d)},
    \qquad
    \mathcal T_{N\to K}^{(d)}\ \text{is entanglement breaking},
    \qquad
    \mathcal C_{N\to N+K}^{(d)}\ \text{is degradable}.
    \label{eq:p53-cloning-family}
  \end{equation}
  Equation~\eqref{eq:p53-cloning-family} eliminates the cloning channels that
  motivated conjugate degradability, but does not prove a general containment
  theorem \sourcecite{ref:p53-brzic-et-al}{BGS+26}.

  \item Zhu and Wang's Theorem~1 and Appendix~A construct a transpose-antidegradable channel $\mathcal N$ whose product with an antidegradable erasure channel has positive private capacity. Thus $\mathcal N$ is not antidegradable, and its complement is a strict example for Eq.~\eqref{eq:p53-transpose-degradable} \sourcecite{ref:p53-zhu-wang-private}{ZW26}.

  \item An independent construction gives $\mathcal F(\rho)=\frac16\sum_{P\in\mathscr S}P\rho P^\dagger$, with $\mathscr S=\{X\otimes X,Y\otimes I,I\otimes Z,X\otimes Z,I\otimes Y,Y\otimes Z\}$. Section~2, Eqs.~(3)--(8), proves transpose antidegradability; Section~4, Eq.~(50), excludes ordinary antidegradability through private-capacity activation. Hence $\Phi=\mathcal F^c$ satisfies the requested separation, with input, output, and environment dimensions $(4,6,4)$ \sourcecite{ref:p53-pauli-private}{GHE26}.
\end{itemize}

\subsection*{References}

\begin{enumerate}[label={},leftmargin=4.8em,labelsep=0.5em]
  \footnotesize
  \item[\textup{[SD22]}]\label{ref:p53-singh-datta}
  S. Singh and N. Datta,
  \lq\lq Detecting Positive Quantum Capacities of Quantum Channels,\rq\rq{}
  \emph{npj Quantum Information} \textbf{8}, 50 (2022).
  \href{https://doi.org/10.1038/s41534-022-00550-2}{doi:10.1038/s41534-022-00550-2};
  \href{https://arxiv.org/abs/2105.06327}{arXiv:2105.06327}.

  \item[\textup{[Bra15]}]\label{ref:p53-bradler}
  K. Br{\'a}dler,
  \lq\lq The Pitfalls of Deciding Whether a Quantum Channel Is (Conjugate)
  Degradable and How to Avoid Them,\rq\rq{}
  \emph{Open Systems \& Information Dynamics} \textbf{22}, 1550026 (2015).
  \href{https://doi.org/10.1142/S1230161215500262}{doi:10.1142/S1230161215500262};
  \href{https://arxiv.org/abs/1507.06159}{arXiv:1507.06159}.

  \item[\textup{[SG16]}]\label{ref:p53-siddhu-griffiths}
  V. Siddhu and R. B. Griffiths,
  \lq\lq Degradable Quantum Channels Using Pure-State to Product-of-Pure-State
  Isometries,\rq\rq{} \emph{Physical Review A} \textbf{94}, 052331 (2016).
  \href{https://doi.org/10.1103/PhysRevA.94.052331}{doi:10.1103/PhysRevA.94.052331};
  \href{https://arxiv.org/abs/1511.05532}{arXiv:1511.05532}.

  \item[\textup{[BGS+26]}]\label{ref:p53-brzic-et-al}
  V. Brzi{\'c}, D. Grinko, M. Studzi{\'n}ski, and M. T. Quintino,
  \lq\lq Optimal Pure State Cloning and Transposition Are Complementary
  Channels,\rq\rq{} arXiv preprint (2026).
  \href{https://arxiv.org/abs/2603.23628}{arXiv:2603.23628}.

  \item[\textup{[ZW26]}]\label{ref:p53-zhu-wang-private}
  C. Zhu and X. Wang, ``Private Communication via Zero-Private-Capacity Quantum Channels,'' arXiv preprint, version 1, 9 September 2026. \href{https://arxiv.org/abs/2609.10520v1}{arXiv:2609.10520v1}.

  \item[\textup{[GHE26]}]\label{ref:p53-pauli-private}
  Uthirakalyani G, P. Halder, and D. Elkouss, ``Private Communication from Pauli Channels with No Privacy,'' arXiv preprint, version 1, 17 September 2026. \href{https://arxiv.org/abs/2609.20133v1}{arXiv:2609.20133v1}.
\end{enumerate}

\subsection*{Comment}

The answer is affirmative. Complementing either cited transpose-antidegradable but non-antidegradable channel gives a channel satisfying Eq.~\eqref{eq:p53-transpose-degradable} but not Eq.~\eqref{eq:p53-degradable}. These complete resolutions are September 2026 preprints, not identified as peer-reviewed at the 18 September 2026 audit. Determining all componentwise-minimal dimensions remains a separate open question.

\subsection*{Field}

Quantum Communication

\subsection*{Topic}

Channel degradability; Quantum channel structure

\subsection*{ID}

\texttt{op\_7e7e4a25fef5c994}
